% ============================================================
% 六、模型的建立与求解
% 高频主线：模型建立 -> 求解方法 -> 结果分析 -> 检验/敏感性
% ============================================================
\section{模型的建立与求解}

\subsection{问题一模型的建立与求解}
\subsubsection{模型准备与变量定义}
\subsubsection{模型建立}
\subsubsection{求解方法与算法流程}
\subsubsection{结果与分析}
\subsubsection{模型检验或敏感性分析}
% 若本问无需单独检验，可删除上一小节。

\subsection{问题二模型的建立与求解}
\subsubsection{模型准备与问题一模型的承接}
\subsubsection{模型建立或修正}
\subsubsection{求解方法与算法流程}
\subsubsection{结果与分析}
\subsubsection{模型检验或敏感性分析}

\subsection{问题三模型的建立与求解}
\subsubsection{模型准备与前述模型的承接}
\subsubsection{模型建立或修正}
\subsubsection{求解方法与算法流程}
\subsubsection{结果与分析}
\subsubsection{对比分析或模型检验}

\subsection{问题四模型的建立与求解}
\subsubsection{模型准备与前述模型的承接}
\subsubsection{模型建立或扩展}
\subsubsection{求解方法与算法流程}
\subsubsection{结果与分析}
\subsubsection{模型检验或讨论}

% 若赛题有第五问，可复制以下结构：
% \subsection{问题五模型的建立与求解}
% \subsubsection{模型建立}
% \subsubsection{求解方法与算法流程}
% \subsubsection{结果与分析}

\subsection{公式、图表与交叉引用示例}
\begin{equation}
 E = mc^2
 \label{eq:placeholder-energy}
\end{equation}
式\eqref{eq:placeholder-energy}为带编号公式。
\begin{figure}[htbp]
\centering
\includegraphics[width=0.55\textwidth]{example_f1.png}
\caption{占位图片示例}
\label{fig:placeholder-example}
\end{figure}
图\ref{fig:placeholder-example}为带编号图片。
\begin{table}[htbp]
\centering
\caption{占位三线表格示例}
\label{tab:placeholder-example}
\begin{tabular}{ccc}
\toprule
指标 & 符号 & 示例值 \\
\midrule
变量一 & $x$ & 1.00 \\
\bottomrule
\end{tabular}
\end{table}
表\ref{tab:placeholder-example}为带编号表格。
